\documentclass{article}

\usepackage{arxiv}

\usepackage[utf8]{inputenc}
\usepackage[T1]{fontenc}
\usepackage{hyperref}
\usepackage{url}
\usepackage{booktabs}
\usepackage{amsfonts}
\usepackage{nicefrac}
\usepackage{microtype}
\usepackage{amsmath}
\usepackage{graphicx}
\usepackage{float}
\graphicspath{ {./graficas_tp21/} }

\title{TP 2.1 - GENERADORES PSEUDOALEATORIOS}

\author{
  Brizio Augusto \\
  Universidad Tecnológica Nacional - FRRO \\
  Zeballos 1341, S2000, Rosario, Argentina \\
  \texttt{brizioaugusto@gmail.com} \\
  \And
  Casermeiro Gonzalo Nicolás \\
  Universidad Tecnológica Nacional - FRRO \\
  Zeballos 1341, S2000, Rosario, Argentina \\
  \texttt{gonzacasermeiro@gmail.com} \\
  \And
  Messina Tiziano Leonel \\
  Universidad Tecnológica Nacional - FRRO \\
  Zeballos 1341, S2000, Rosario, Argentina \\
  \texttt{tizilmessina@gmail.com} \\
  \And
  De la Rosa Valentin \\
  Universidad Tecnológica Nacional - FRRO \\
  Zeballos 1341, S2000, Rosario, Argentina \\
  \texttt{mferreyra079@gmail.com} \\
}

\begin{document}
\maketitle

\begin{abstract}
Este trabajo estudia la generación de números pseudoaleatorios mediante la
implementación y comparación de tres fuentes de datos: un generador congruencial
lineal (GCL), el método de los cuadrados medios y el generador provisto por
Python mediante \texttt{random.random()}. Para evaluar la calidad estadística de
las secuencias generadas se aplicaron seis pruebas: test de media, test de
varianza, prueba chi-cuadrado de uniformidad, prueba de Kolmogorov-Smirnov,
prueba de corridas y autocorrelación de primer orden.

Se generaron $n = 10\,000$ valores para cada generador utilizando una semilla
fija para garantizar reproducibilidad. Los resultados muestran que el GCL
implementado y el generador de Python superan todas las pruebas consideradas,
mientras que el método de cuadrados medios presenta fallas significativas en
uniformidad, media y varianza. El análisis confirma que no alcanza con producir
valores aparentemente desordenados: un generador útil para simulación debe
aproximar una distribución uniforme, evitar patrones detectables y sostener esas
propiedades a lo largo de secuencias suficientemente largas.
\end{abstract}

\section{Introducción}

La simulación computacional requiere representar procesos inciertos mediante
secuencias numéricas que se comporten como si fueran aleatorias. En una
computadora, sin embargo, la aleatoriedad pura no surge naturalmente: los valores
son generados por algoritmos deterministas que, a partir de un estado inicial o
\textit{semilla}, producen una secuencia reproducible. A estos algoritmos se los
conoce como generadores de números pseudoaleatorios.

La reproducibilidad es una ventaja central en simulación, ya que permite repetir
un experimento bajo las mismas condiciones iniciales y verificar los resultados.
Pero también introduce un riesgo: si el algoritmo posee ciclos cortos, sesgos o
correlaciones internas, la simulación puede producir conclusiones equivocadas.
Por este motivo, además de programar generadores, es necesario evaluarlos con
pruebas estadísticas que permitan detectar desviaciones respecto del
comportamiento esperado.

El objetivo de este trabajo es implementar generadores pseudoaleatorios en
Python 3.x, aplicar tests de calidad sobre las secuencias obtenidas y comparar
su desempeño con el generador incluido en el lenguaje. El informe se organiza de
la siguiente manera: la Sección~2 describe los generadores estudiados; la
Sección~3 presenta las pruebas estadísticas utilizadas; la Sección~4 expone los
resultados obtenidos; y la Sección~5 desarrolla las conclusiones.

\section{Generadores estudiados}

\subsection{Generador congruencial lineal}

El generador congruencial lineal (GCL) es uno de los métodos clásicos para
obtener números pseudoaleatorios. A partir de una semilla inicial $X_0$, la
secuencia se define mediante la recurrencia:

\begin{equation}
    X_{i+1} = (aX_i + c) \bmod m
\end{equation}

\noindent donde $a$ es el multiplicador, $c$ el incremento y $m$ el módulo. Para
trabajar con valores en el intervalo $[0,1)$, cada entero generado se normaliza:

\begin{equation}
    U_i = \frac{X_i}{m}
\end{equation}

En la simulación se utilizaron los parámetros:

\begin{equation}
    a = 1103515245, \qquad c = 12345, \qquad m = 2^{31}
\end{equation}

\noindent con semilla $X_0 = 5731$. Estos parámetros corresponden a una familia
de GCLs ampliamente utilizada históricamente. Si bien los GCLs modernos no son
recomendables para criptografía, siguen siendo útiles para introducir los
principios de generación pseudoaleatoria y para comparar calidad estadística.

\subsection{Método de los cuadrados medios}

El método de los cuadrados medios fue propuesto por John von Neumann. La idea es
elevar al cuadrado la semilla actual, completar con ceros si es necesario y tomar
los dígitos centrales como nuevo estado:

\begin{equation}
    X_{i+1} = \text{dígitos centrales}(X_i^2)
\end{equation}

Luego se normaliza el valor dividiéndolo por $10^d$, donde $d$ es la cantidad de
dígitos utilizada:

\begin{equation}
    U_i = \frac{X_i}{10^d}
\end{equation}

En este trabajo se emplearon semillas de cuatro dígitos con $d = 4$ y semilla inicial
$5731$. Aunque el método es simple, suele presentar ciclos cortos, convergencia a
cero o patrones repetitivos, por lo que se espera un rendimiento inferior al del
GCL y al generador de Python.

\subsection{Generador de Python}

Como referencia se utilizó \texttt{random.random()} de Python, que genera valores
en $[0,1)$ mediante el generador Mersenne Twister. Este generador tiene un período
muy grande y buenas propiedades estadísticas para simulación general, aunque no
debe utilizarse para fines criptográficos \cite{python_random}.

\section{Pruebas estadísticas}

Para cada generador se produjeron $n = 10\,000$ valores pseudoaleatorios. Las
secuencias fueron sometidas a seis pruebas. En todas ellas se tomó un nivel de
significancia aproximado de $\alpha = 0{,}05$, equivalente al criterio usual
$|Z| < 1{,}96$ para pruebas bilaterales.

\subsection{Test de media}

Para una distribución uniforme $U(0,1)$, la esperanza teórica es:

\begin{equation}
    E[U] = \frac{1}{2}
\end{equation}

El estadístico utilizado fue:

\begin{equation}
    Z_{\bar{U}} =
    \frac{|\bar{U} - 0{,}5|}{\sqrt{\frac{1}{12n}}}
\end{equation}

\noindent y se acepta el comportamiento esperado si $Z_{\bar{U}} < 1{,}96$.

\subsection{Test de varianza}

La varianza teórica de una distribución uniforme $U(0,1)$ es:

\begin{equation}
    Var(U) = \frac{1}{12}
\end{equation}

Se comparó la varianza muestral contra este valor esperado mediante una
aproximación normal del estadístico chi-cuadrado:

\begin{equation}
    Z_{S^2} =
    \frac{\left| \frac{(n-1)S^2}{1/12} - (n-1) \right|}
    {\sqrt{2(n-1)}}
\end{equation}

\subsection{Prueba chi-cuadrado}

El intervalo $[0,1)$ se dividió en $k = 10$ subintervalos de igual amplitud. Para
cada intervalo se comparó la frecuencia observada $O_j$ contra la frecuencia
esperada $E_j = n/k$:

\begin{equation}
    \chi^2 = \sum_{j=1}^{k} \frac{(O_j - E_j)^2}{E_j}
\end{equation}

Con $k-1 = 9$ grados de libertad, el valor crítico para $\alpha = 0{,}05$ es
$16{,}919$. La secuencia supera la prueba si $\chi^2 < 16{,}919$.

\subsection{Prueba de Kolmogorov-Smirnov}

La prueba de Kolmogorov-Smirnov compara la distribución empírica $F_n(x)$ con la
función de distribución acumulada teórica $F(x) = x$:

\begin{equation}
    D = \sup_x |F_n(x) - F(x)|
\end{equation}

Para $n = 10\,000$, se utilizó el criterio:

\begin{equation}
    D < \frac{1{,}36}{\sqrt{n}} = 0{,}01360
\end{equation}

\subsection{Prueba de corridas}

La prueba de corridas analiza la alternancia de valores por encima y por debajo
de $0{,}5$. Si los valores son independientes, la cantidad de corridas debe
aproximarse a:

\begin{equation}
    E[R] = \frac{2n_1n_2}{n_1+n_2} + 1
\end{equation}

\noindent donde $n_1$ es la cantidad de valores superiores o iguales a $0{,}5$ y
$n_2$ la cantidad de valores inferiores.

\subsection{Autocorrelación}

Para detectar dependencia lineal entre valores consecutivos se calculó la
autocorrelación de primer orden:

\begin{equation}
    \rho_1 =
    \frac{\sum_{i=1}^{n-1}(U_i-\bar{U})(U_{i+1}-\bar{U})}
    {\sum_{i=1}^{n}(U_i-\bar{U})^2}
\end{equation}

Luego se evaluó el estadístico $Z = |\rho_1|\sqrt{n}$, aceptando la hipótesis de
ausencia de autocorrelación lineal si $Z < 1{,}96$.

\section{Resultados y discusión}

La Tabla~\ref{tab:resultados} resume los resultados obtenidos para los tres
generadores. El valor \textbf{OK} indica que el estadístico quedó dentro del
rango esperado; \textbf{RECHAZA} indica evidencia suficiente para descartar el
comportamiento esperado en esa prueba.

\begin{table}[H]
\centering
\caption{Resultados de las pruebas estadísticas para $n = 10\,000$ valores.}
\label{tab:resultados}
\begin{tabular}{llcl}
\toprule
\textbf{Generador} & \textbf{Test} & \textbf{Estadístico} & \textbf{Resultado} \\
\midrule
GCL & Media & 0{,}04387 & OK \\
GCL & Varianza & 0{,}01743 & OK \\
GCL & Chi-cuadrado & 3{,}14800 & OK \\
GCL & Kolmogorov-Smirnov & 0{,}00492 & OK \\
GCL & Corridas & 0{,}99737 & OK \\
GCL & Autocorrelación & 1{,}25931 & OK \\
\midrule
Cuadrados medios & Media & 3{,}32149 & RECHAZA \\
Cuadrados medios & Varianza & 27{,}99503 & RECHAZA \\
Cuadrados medios & Chi-cuadrado & 14795{,}75200 & RECHAZA \\
Cuadrados medios & Kolmogorov-Smirnov & 0{,}20780 & RECHAZA \\
Cuadrados medios & Corridas & 0{,}05990 & OK \\
Cuadrados medios & Autocorrelación & 0{,}18212 & OK \\
\midrule
Python random & Media & 0{,}88573 & OK \\
Python random & Varianza & 0{,}23982 & OK \\
Python random & Chi-cuadrado & 4{,}69600 & OK \\
Python random & Kolmogorov-Smirnov & 0{,}00903 & OK \\
Python random & Corridas & 0{,}43448 & OK \\
Python random & Autocorrelación & 0{,}12095 & OK \\
\bottomrule
\end{tabular}
\end{table}

\begin{figure}[H]
    \centering
    \includegraphics[width=\linewidth]{histograma_generadores.png}
    \caption{Histograma comparativo de los valores generados por cada método.}
    \label{fig:histograma}
\end{figure}

\begin{figure}[H]
    \centering
    \includegraphics[width=\linewidth]{primeros_valores.png}
    \caption{Primeros 200 valores generados por cada método.}
    \label{fig:primeros_valores}
\end{figure}

\begin{figure}[H]
    \centering
    \includegraphics[width=\linewidth]{dispersion_pares.png}
    \caption{Diagrama de dispersión entre valores consecutivos $U_i$ y $U_{i+1}$.}
    \label{fig:dispersion}
\end{figure}

\subsection{Análisis del GCL}

El generador congruencial lineal superó las seis pruebas aplicadas. La media y
la varianza se mantuvieron muy cerca de los valores teóricos de una uniforme
$U(0,1)$, mientras que las pruebas de uniformidad global también quedaron dentro
de los límites esperados. El valor de chi-cuadrado fue $3{,}14800$, muy por
debajo del límite $16{,}919$, lo que indica una distribución equilibrada entre
los subintervalos.

La prueba de corridas y la autocorrelación tampoco detectaron patrones fuertes
entre valores consecutivos. Esto no significa que el GCL sea perfecto: los GCLs
pueden presentar estructuras en dimensiones superiores y no son aptos para
criptografía. Sin embargo, para el tamaño de muestra y las pruebas usadas en
este trabajo, su comportamiento fue aceptable para simulación básica.

\subsection{Análisis del método de cuadrados medios}

El método de cuadrados medios mostró el peor desempeño. Falló el test de media,
el test de varianza, la prueba chi-cuadrado y la prueba de Kolmogorov-Smirnov.
El rechazo por chi-cuadrado fue especialmente marcado, con un
$14795{,}75200$, muy superior al valor crítico $16{,}919$. Esto indica que los
valores no se distribuyeron uniformemente en el intervalo $[0,1)$.

Aunque superó la prueba de corridas y la prueba de autocorrelación, esto no
alcanza para considerarlo un generador adecuado. Una secuencia puede alternar
razonablemente alrededor de $0{,}5$ y aun así estar concentrada en ciertas zonas
del intervalo o tener una varianza incorrecta. El resultado confirma la debilidad
histórica del método: es simple de implementar, pero muy sensible a la semilla y
propenso a ciclos cortos o estados degenerados.

\subsection{Análisis del generador de Python}

El generador de Python superó todas las pruebas, como era esperable para una
herramienta estándar de simulación general. Sus estadísticos se mantuvieron lejos
de los valores críticos tanto en uniformidad como en independencia básica. En
particular, la prueba chi-cuadrado obtuvo $4{,}69600$ y Kolmogorov-Smirnov
$0{,}00903$, valores compatibles con una muestra uniforme.

Este comportamiento lo convierte en una alternativa confiable para simulaciones
académicas y experimentales. No obstante, como indica la documentación oficial,
el módulo \texttt{random} no debe utilizarse para aplicaciones criptográficas,
ya que su objetivo es la simulación reproducible y no la seguridad.

\section{Conclusiones}

El trabajo permitió comprobar experimentalmente que la calidad de un generador
pseudoaleatorio no puede asumirse únicamente porque sus valores parezcan
desordenados. Es necesario aplicar pruebas estadísticas que evalúen uniformidad,
media, varianza, independencia y ausencia de patrones simples.

El generador congruencial lineal implementado tuvo un desempeño satisfactorio en
las seis pruebas realizadas. Esto muestra que, con una selección adecuada de
parámetros, un GCL puede ser útil en simulaciones introductorias. Sin embargo,
también se debe reconocer su limitación: existen pruebas más exigentes y
análisis multidimensionales capaces de revelar patrones que los tests básicos no
detectan.

El método de los cuadrados medios resultó claramente deficiente. Su falla en
media, varianza y uniformidad confirma que se trata de un algoritmo histórico y
didáctico, pero poco recomendable para simulaciones serias. La comparación con
Python refuerza esta conclusión: el generador estándar del lenguaje mantuvo un
comportamiento estable y superó todas las pruebas propuestas.

En síntesis, para simulaciones reales conviene utilizar generadores probados y
bien documentados. Los métodos simples como cuadrados medios son valiosos para
comprender el problema, pero no para reemplazar generadores modernos. El GCL,
por su parte, funciona correctamente en este experimento y permite observar de
manera clara la relación entre algoritmo, semilla, período y calidad estadística.

\bibliographystyle{unsrt}

\begin{thebibliography}{9}

\bibitem{torres_tp21}
J.~I. Torres, \textit{TP 2.1 - Generadores pseudoaleatorios},
Universidad Tecnológica Nacional - FRRO, 2026.

\bibitem{tereom}
Tereom, \textit{Números pseudoaleatorios},
\url{https://tereom.github.io/est-computacional-2018/numeros-pseudoaleatorios.html}.

\bibitem{random_analysis}
Random.org, \textit{Introduction to Randomness and Random Numbers},
\url{https://www.random.org/analysis/}.

\bibitem{random_analysis_pdf}
Random.org, \textit{Analysis of the Random Number Generator},
\url{https://www.random.org/analysis/Analysis2005.pdf}.

\bibitem{python_random}
Python Software Foundation, \textit{random --- Generate pseudo-random numbers},
\url{https://docs.python.org/3/library/random.html}.

\bibitem{knuth}
D.~E. Knuth, \textit{The Art of Computer Programming, Volume 2:
Seminumerical Algorithms}, Addison-Wesley, 1997.

\end{thebibliography}

\end{document}
